\documentclass[Main.tex]{subfiles}
\begin{document}
\section{Ôn tập Chương 3}
\subsection{Lý thuyết trọng tâm}
\begin{kienthuccannho}
	\subsubsection{Vật liệu}
	\begin{itemize}
		\item \textbf{\textit{Vật liệu}} là chất hoặc hỗn hợp một số chất được con người sử dụng như nguyên liệu đầu vào trong quá trình sản xuất hoặc chế tạo để làm ra những sản phẩm phục vụ cuộc sống.
		\item Tính chất và ứng dụng của một số vật liệu:
		\begin{itemize}
			\item \textbf{Kim loại:} dẻo, dẫn điện, dẫn nhiệt tốt, có ánh kim, bị ăn mòn. Làm dây điện, xoong nồi, …
			\item \textbf{Thủy tinh:} trong suốt, cho ánh sáng đi qua, giòn, dễ vỡ. Làm bình hoa, chai lọ, cửa kính, …
			\item \textbf{Nhựa:} dẻo, nhẹ, không dẫn điện, dễ bị biến dạng nhiệt. Làm ghế ngồi, ống dẫn nước, …
			\item \textbf{Gốm, sứ:} không bị ăn mòn, dẫn nhiệt kém, giòn, dễ vỡ. Làm chum vại, bát đĩa, chậu hoa, …
			\item \textbf{Cao su:} đàn hồi, bền, không thấm nước. Làm lốp xe, gioăng cao su, đệm, …
			\item \textbf{Gỗ:} bền, chịu lực tốt, dễ cháy, có thể bị mối mọt. Làm nhà, khung cửa, bàn, ghế, tủ, …
		\end{itemize}
	\end{itemize}
	\subsubsection{Nhiên liệu}
	\begin{itemize}
		\item \textbf{\textit{Nhiên liệu}} (chất đốt) là những chất cháy được và khi cháy tỏa nhiều nhiệt.
		\item Tính chất và ứng dụng của một số nhiên liệu:
		\begin{itemize}
			\item \textbf{Than đá:} chất rắn, cháy tỏa nhiều nhiệt. Dùng để đun nấu, sưởi ấm, nhiên liệu công nghiệp, …
			\item \textbf{Xăng dầu:} chất lỏng, cháy tỏa nhiều nhiệt. Dùng để đun nấu, chạy động cơ, …
			\item \textbf{Khí gas:} chất khí, cháy tỏa nhiều nhiệt. Dùng trong đun nấu, …
		\end{itemize}
		\item \textbf{\textit{An ninh năng lượng}} là sự đảm bảo đầy đủ năng lượng dưới nhiều dạng khác nhau, ưu tiên các nguồn năng lượng sạch và giá thành rẻ.
	\end{itemize}
	\subsubsection{Nguyên liệu}
	\begin{itemize}
		\item \textbf{\textit{Nguyên liệu}} là vật liệu tự nhiên (vật liệu thô) chưa qua xử lí và cần chuyển hóa để tạo ra sản phẩm.
		\item Tính chất và ứng dụng của một số nguyên liệu:
		\begin{itemize}
			\item \textbf{Đá vôi:} thành phần chính là calcium carbonate, có nhiều màu, cứng, không tan trong nước, tan trong acid. Dùng để sản xuất vôi sống, xây dựng, làm chất độn trong sản xuất cao su, xà phòng, …
			\item \textbf{Quặng:} có màu sắc đa dạng. Quặng bauxite sản xuất nhôm, quặng sắt sản xuất gang thép.
		\end{itemize}
	\end{itemize}
	\subsubsection{Lương thực – thực phẩm}
	\begin{itemize}
		\item \textbf{\textit{Lương thực}} như gạo, ngô, khoai, sắn, … chứa nhiều tinh bột. \textbf{\textit{Thực phẩm}} như thịt, cá, trứng, sữa, … dùng làm các món ăn.
		\item Thành phần chất dinh dưỡng trong lương thực -- thực phẩm:
		\begin{itemize}
			\item \textbf{Carbohydrate:} chứa nhiều tinh bột, đường, chất xơ, là nguồn cung cấp năng lượng chính cho cơ thể hoạt động. Có nhiều trong lúa, ngô, khoai, sắn, …
			\item \textbf{Protein (đạm):} liên quan đến mọi chức năng sống của cơ thể. Có nhiều trong cá, thịt, trứng, sữa, đậu đỗ, …
			\item \textbf{Lipid (chất béo):} là nguồn năng lượng dự trữ trong cơ thể. Có nhiều trong sữa, thịt, cá, lạc, …
			\item \textbf{Vitamin và khoáng chất:} nâng cao hệ miễn dịch của cơ thể. Có nhiều trong hải sản, rau xanh, củ quả tươi, …
		\end{itemize}
		\item Tất cả vật liệu, nguyên liệu, nhiên liệu và lương thực -- thực phẩm phải được sử dụng tiết kiệm, an toàn và đảm bảo sự phát triển bền vững.
	\end{itemize}
\end{kienthuccannho}
\subsection{Bài tập}

%%%=================Phần bài tập tự luận===================%%%
\phan{Bài tập tự luận}
\textit{\large Thí sinh trả lời từ câu 1 đến câu 6}
\Opensolutionfile{ansbth}[Ans/LGBT-C03_ON_TAP_CHUONG]
\Opensolutionfile{ansbt}[Ans/AnsBT-C03_ON_TAP_CHUONG]

%%%=============BT_1=============%%%
\begin{bt}%[6K3H1-1]
	Liệt kê các từ (in nghiêng) chỉ vật liệu và nguyên liệu trong các trường hợp sau đây:
	\begin{enumerate}
		\item[(a)] \textit{Nhựa} dùng để làm bàn, ghế, cốc, tấm lợp, …
		\item[(b)] \textit{Gỗ} dùng để làm bàn, ghế, nhà cửa, …
		\item[(c)] \textit{Đá vôi} được xay vụn để trải đường, trộn bê tông, …
		\item[(d)] \textit{Đá vôi} được nung lên để lấy sản phẩm xây nhà, khử chua cho ruộng đồng, …
		\item[(e)] \textit{Quặng sắt} được nung trong lò cao tạo ra gang, thép để làm nhà, cửa, cầu cống, ống dẫn nước, các bộ phận máy móc, …
	\end{enumerate}
	\loigiai{
		\begin{itemize}
			\item \textbf{Vật liệu:} nhựa, gỗ, đá vôi ở trường hợp (c).
			\item \textbf{Nguyên liệu:} đá vôi ở trường hợp (d), quặng sắt.
		\end{itemize}
	}
\end{bt}

%%%=============BT_2=============%%%
\begin{bt}%[6K3V1-6]
	Ghép mỗi loại rác thải với cách xử lí phù hợp:
	\begin{center}
		\renewcommand{\arraystretch}{1.15}
		\begin{tabular}{|L{0.28\linewidth}|L{0.58\linewidth}|}
			\hline
			\textbf{Loại rác thải} & \textbf{Cách xử lí} \\
			\hline
			$(1)$ Thức ăn thừa, rác thực vật, xác động vật & (a) là nguyên liệu dùng làm phân bón bằng cách băm nhỏ và cho vào một thùng xốp để phân hủy hoàn toàn (có thể sử dụng thùng xốp để trồng rau sạch) hoặc đào hố để ủ rác làm phân bón. \\
			\hline
			$(2)$ Đồ điện hư hỏng & (b) có nhiều chất độc hại cần phải được xử lí chuyên nghiệp để tái chế vào việc khác hoặc thu hồi các thành phần của nó một cách an toàn. \\
			\hline
			$(3)$ Pin điện, nguồn điện hỏng & (c) từng bộ phận riêng biệt có chức năng và tính chất khác nhau nên cần được chuyển đến đúng nơi có đủ chức năng xử lí theo đúng quy định, đảm bảo an toàn. \\
			\hline
		\end{tabular}
	\end{center}
	\loigiai{
		Kết quả ghép: $(1)$ -- (a), $(2)$ -- (c), $(3)$ -- (b).
	}
\end{bt}

%%%=============BT_3=============%%%
\begin{bt}%[6K3H1-1]
	Chọn những vật liệu phù hợp (kim loại, thủy tinh, nhựa, gốm sứ, cao su, gỗ) để hoàn thành bảng thông tin sau:
	\begin{center}
		\renewcommand{\arraystretch}{1.15}
		\begin{tabular}{|L{0.16\linewidth}|L{0.42\linewidth}|L{0.28\linewidth}|}
			\hline
			\textbf{Vật liệu} & \textbf{Tính chất} & \textbf{Ứng dụng} \\
			\hline
			$(1)$ ……….. & Có ánh kim, dẫn điện tốt, dẫn nhiệt tốt, dẻo có thể kéo thành sợi và dát mỏng, có thể bị gỉ. & Làm dây dẫn điện, nồi đun nấu, làm cầu cống, khung nhà cửa, … \\
			\hline
			$(2)$ ……….. & Trong suốt, cho ánh sáng đi qua, dẫn nhiệt kém, không dẫn điện, cứng nhưng giòn, dễ vỡ. & Làm bình hoa, chai lọ, cửa kính, … \\
			\hline
			$(3)$ ……….. & Dẻo, nhẹ, không dẫn điện, dẫn nhiệt kém, không bị ăn mòn, dễ bị biến dạng nhiệt. & Làm ghế ngồi, ống dẫn nước, tấm lợp, … \\
			\hline
			$(4)$ ……….. & Không bị ăn mòn, dẫn nhiệt kém, hầu như không dẫn điện, cứng nhưng giòn, dễ vỡ. & Làm chum vại, bát đĩa, chậu hoa, … \\
			\hline
			$(5)$ ……….. & Đàn hồi, bền, không dẫn điện và nhiệt, không thấm nước, dễ cháy. & Làm lốp xe, gioăng cao su, đệm, … \\
			\hline
			$(6)$ ……….. & Bền, chịu lực tốt, dễ tạo hình, dễ cháy, có thể bị mối mọt. & Làm nhà, khung cửa, bàn, ghế, tủ, … \\
			\hline
		\end{tabular}
	\end{center}
	\loigiai{
		\begin{itemize}
			\item $(1)$ Kim loại.
			\item $(2)$ Thủy tinh.
			\item $(3)$ Nhựa.
			\item $(4)$ Gốm, sứ.
			\item $(5)$ Cao su.
			\item $(6)$ Gỗ.
		\end{itemize}
	}
\end{bt}

%%%=============BT_4=============%%%
\begin{bt}%[6K3V4-1]
	Cơ thể con người cần nhiều loại chất khác nhau và hầu hết các chất này đều lấy từ thức ăn. Hãy điền những chất dinh dưỡng chính tương ứng với mỗi loại thức ăn sau:
	\begin{center}
		\renewcommand{\arraystretch}{1.15}
		\begin{tabular}{|L{0.35\linewidth}|L{0.5\linewidth}|}
			\hline
			\textbf{Thức ăn} & \textbf{Chất dinh dưỡng} \\
			\hline
			Cá basa kho tộ & \\
			\hline
			Canh chua cá lóc & \\
			\hline
			Gà luộc & \\
			\hline
			Rau cải xanh xào tỏi & \\
			\hline
			Cơm trắng & \\
			\hline
		\end{tabular}
	\end{center}
	\loigiai{
		\begin{center}
			\renewcommand{\arraystretch}{1.15}
			\begin{tabular}{|L{0.35\linewidth}|L{0.5\linewidth}|}
				\hline
				\textbf{Thức ăn} & \textbf{Chất dinh dưỡng} \\
				\hline
				Cá basa kho tộ & Protein, lipid, calcium và một số vi chất khác. \\
				\hline
				Canh chua cá lóc & Protein, lipid, calcium và vitamin (từ quả chua như dứa, cà chua, …). \\
				\hline
				Gà luộc & Protein, lipid. \\
				\hline
				Rau cải xanh xào tỏi & Vitamin và một số chất khoáng. \\
				\hline
				Cơm trắng & Tinh bột. \\
				\hline
			\end{tabular}
		\end{center}
	}
\end{bt}

%%%=============BT_5=============%%%
\begin{bt}%[6K3V3-3]
	Ủy ban nhân dân thành phố Hà Nội đã ra quyết định ngừng sử dụng bếp than tổ ong ở các quận nội thành Hà Nội bắt đầu từ ngày $1 - 1 - 2021$. Em hãy cho biết lí do vì sao?
	\loigiai{
		\begin{itemize}
			\item Nhiên liệu hóa thạch thường chứa nhiều tạp chất, khi đốt sẽ hình thành nhiều chất có hại đến môi trường.
			\item Ngay cả khi chỉ có carbon hay các hydrocarbon thì việc đốt cháy không hoàn toàn cũng sinh ra khí độc là carbon monoxide (CO).
			\item Khi làm than tổ ong, người ta thường trộn với bùn nên lại có thêm tạp chất, có thể gây ô nhiễm hơn nữa.
		\end{itemize}
		Vì vậy, để giữ cho thành phố Hà Nội sạch đẹp và an toàn thì việc dừng sử dụng than tổ ong là hợp lí.
	}
\end{bt}

%%%=============BT_6=============%%%
\begin{bt}%[6K3V1-1]
	Em hãy tìm hiểu qua sách, báo, internet, … và kể một số vật liệu, nguyên liệu, nhiên liệu và lương thực -- thực phẩm có ở nước ta.
	\loigiai{
		\begin{itemize}
			\item \textbf{Vật liệu tự nhiên:} gỗ, tre, đá vôi, cát, sỏi, …
			\item \textbf{Nguyên liệu lấy từ tự nhiên:} đá vôi, quặng sắt, bauxite, đất sét, gỗ, …
			\item \textbf{Nhiên liệu lấy từ tự nhiên:} than đá, dầu mỏ, khí thiên nhiên, …
			\item \textbf{Lương thực -- thực phẩm:} vô cùng đa dạng và phong phú ở các vùng miền khác nhau như ngũ cốc (gạo, lúa mì, ngô, khoai, sắn), thịt gia súc (trâu, bò, …), thịt gia cầm (gà, vịt, ngan, ngỗng, chim, …) và các sản phẩm từ gia cầm (trứng), …
		\end{itemize}
	}
\end{bt}
\Closesolutionfile{ansbt}
\Closesolutionfile{ansbth}
\end{document}
